% Stage 4 (author: Himanshu Kumar).
% Training numbers come from results/stage4_masked_densenet/summary.json and
% results/stage2_phaseA/settings.json / results/stage2_phaseB/lr1e-4_bs16/settings.json.
% Test numbers come from results/stage4_metrics.csv,
% results/stage4_cross_test_metrics.csv, results/stage4_comparison.csv,
% results/stage4_masked_densenet/masked_test_predictions_attention.csv,
% results/stage2_densenet/test_predictions.csv and results/stage3/lung_mask_summary.csv.
\section{Lung-Masked Training and Cross-Test}
\label{sec:masking}

\subsection{Training protocol}
We train a second DenseNet-121 after masking pixels outside the lungs. Each X-ray is converted to grayscale,
resized to $224\times224$, and set to black outside the binary lung mask. The same committed patient-level
splits and Stage~2 training settings are used. Phase~A trains the classifier head and Phase~B fine-tunes the
whole network; the checkpoint is selected by validation loss, without evaluating the test set.

The selected masked model reached its lowest Phase~A validation loss of 0.083285 at epoch~19
(validation AUC 0.986409), then its lowest Phase~B validation loss of 0.028584 at epoch~10
(validation AUC 0.998512). As in Stage~2, Phase~A used batch size 32 and learning rate $10^{-3}$;
Phase~B used batch size 16 and learning rate $10^{-4}$. Both phases used seed 42 and the same
threshold and preprocessing settings as Stage~2.

\subsection{Test evaluation}
After selection, we evaluate the masked model on masked test images and cross-test the Stage~2 unmasked
checkpoint on those same masked images. For the masked model, the same forward/backward pass supplies both
predictions and Grad-CAM maps. We report the share of heatmap energy inside the lung masks with
\texttt{attention\_in\_lung}, shown as a percentage, alongside accuracy, recall, F1 and AUC.

\begin{table}[t]
  \centering
  \small
  \begin{tabular}{llrrrrr}
    \toprule
    Model & Training data & Acc. & Rec. & F1 & AUC & Lung attention (\%) \\
    \midrule
    Custom CNN   & Unmasked & 0.8013 & 0.9821 & 0.8607 & 0.9016 & 43.41\% \\
    DenseNet-121 & Unmasked & 0.8734 & 0.9974 & 0.9078 & 0.9713 & 39.52\% \\
    DenseNet-121 & Lung-masked & 0.8702 & 0.9949 & 0.9055 & 0.9759 & 51.79\% \\
    \bottomrule
  \end{tabular}
  \caption{Final test comparison on the committed 624-image test split. Baseline and unmasked DenseNet-121 values are from Stage~1 and Stage~3; the masked-model values are from the Stage~4.2/4.3 Kaggle evaluation.}
  \label{tab:masking-comparison}
\end{table}

\paragraph{Cross-test.}
The Stage~2 unmasked checkpoint evaluated on masked test images achieves 0.6314 accuracy, 0.6290 precision,
1.0000 recall, 0.7723 F1 and 0.8625 AUC. It predicts 230 of the 234 NORMAL images as PNEUMONIA
(TN=4, FP=230, FN=0, TP=390). This is a substantial drop from its matched unmasked test result
(0.8734 accuracy and 0.9713 AUC), showing sensitivity to removal of extra-pulmonary image context.
The masked inputs also introduce black backgrounds that the unmasked model did not see during training,
so this cross-test drop may partly reflect an input-distribution change.

\paragraph{Interpretation and limitations.}
On matched inputs, the masked model is close to the unmasked DenseNet-121 in accuracy (0.8702 vs.~0.8734),
recall (0.9949 vs.~0.9974) and F1 (0.9055 vs.~0.9078), with a slightly higher AUC (0.9759 vs.~0.9713).
Its mean lung attention is 51.79\%, compared with 39.52\% for the unmasked model and a mean lung-area
baseline of 37.70\%; the masked model's mean attention-to-area lift is 1.40. The masked model has 79 false
positives versus 78 for the unmasked model, with 68 images in common. Both highlighted Stage~3 examples
(NORMAL2-IM-0252 and NORMAL2-IM-0221) remain false positives after masking. Their corner text shows where
the unmasked model attends, but does not explain most errors; the test-set gap remains. The cross-test
indicates sensitivity to changing pixels outside the lungs, but does not identify a specific shortcut or
establish causation. The Stage~3 segmentation caveat also applies: 19 test masks and 182 training masks are
flagged for possible segmentation failure, and all remain included. The eye check also found partial mask
misses that were not flagged. Grad-CAM attention should therefore be interpreted alongside mask quality
rather than as a definitive measure of causal model reliance.
